\documentclass[aps,prb,twocolumn,floatfix]{revtex4-2}

\usepackage{amsmath,amssymb,amsfonts}
\usepackage{bm}
\usepackage{graphicx}
\usepackage{hyperref}
\usepackage{xcolor}

\graphicspath{{../figures/}}

\newcommand{\Hflat}{H_{\rm flat}}
\newcommand{\Pocc}{P_{\rm occ}}
\newcommand{\Cmat}{C}
\newcommand{\Ns}{N_s}
\newcommand{\Ny}{N_y}
\newcommand{\Nx}{N_x}
\newcommand{\qorb}{q}
\newcommand{\Tr}{\operatorname{Tr}}
\newcommand{\diag}{\operatorname{diag}}
\newcommand{\Heff}{h_{\rm ent}}
\newcommand{\ceff}{c_{\rm eff}}

\begin{document}

\title{Gaussian Entanglement Renormalization of the Flattened Domain-Wall Hamiltonian}

\date{\today}

\begin{abstract}
We describe a quasi-one-dimensional Gaussian entanglement renormalization
procedure for the half-filled ground state of the overcomplete-Wannier flattened
Hamiltonian in the domain-wall Chern-insulator geometry.  The two-dimensional
system is regarded as a one-dimensional chain in the \(y\) direction, with each
row carrying \(q=2N_x\) complex-fermion orbitals.  A binary RG layer groups
neighboring rows, diagonalizes the restricted correlation matrix in each
two-row block, keeps the natural orbitals with largest \(n(1-n)\), and projects
the global Gaussian state into the retained subspace.  This construction is
not a scale-invariant MERA: no disentangler is optimized and the block
isometries are allowed to depend on layer and position.  Its purpose is instead
a smoke test for the MERA hypothesis.  If the domain-wall flattened state has
the expected infrared structure, local Gaussian mode sorting should polarize
gapped bulk modes while retaining mixed modes localized near the two domain
walls, and the full-\(x\) strip entropy should remain compatible with
\(c_{\rm eff}=1\).
\end{abstract}

\maketitle

\section{Flattened Hamiltonian target state}

The microscopic single-particle Hilbert space has dimension
\begin{equation}
  N_{\rm layer}=2N_xN_y,
\end{equation}
where the factor of two is the internal Chern-insulator orbital index.  We use
the same domain-wall benchmark geometry as in the flattened-Hamiltonian
analysis: \(N_x=20\), \(N_y=40\), \(\alpha_1=1\), \(\alpha_2=30\),
\(\mathrm{DW}=\mathrm{True}\), and DW-sector truncation enabled.  The
overcomplete-Wannier spinors are denoted
\begin{equation}
  W_{A+},\quad W_{B+},\quad W_{A-},\quad W_{B-},
\end{equation}
where \(A,B\) label the two trial spinors and \(\pm\) labels the Chern-insulator
band projector used in constructing the Wannier orbital.  After reshaping each
family into a matrix whose columns are Wannier orbitals, the flattened
single-particle Hamiltonian is
\begin{equation}
  \Hflat =
  W_{A+}W_{A+}^{\dagger}
  +W_{B+}W_{B+}^{\dagger}
  -W_{A-}W_{A-}^{\dagger}
  -W_{B-}W_{B-}^{\dagger}.
  \label{eq:hflat}
\end{equation}
In the notebook this matrix is explicitly Hermitian symmetrized,
\begin{equation}
  \Hflat \leftarrow \frac{1}{2}\left(\Hflat+\Hflat^\dagger\right),
\end{equation}
before diagonalization.  Let
\begin{equation}
  \Hflat u_a = \epsilon_a u_a,\qquad
  \epsilon_1\leq\epsilon_2\leq\cdots\leq\epsilon_{N_{\rm layer}}.
\end{equation}
The half-filled Gaussian ground state is obtained by occupying
\begin{equation}
  N_{\rm occ}=N_{\rm layer}/2
\end{equation}
single-particle modes.  Equivalently, its occupied-band projector is
\begin{equation}
  \Pocc = \sum_{a=1}^{N_{\rm occ}} u_a u_a^\dagger.
  \label{eq:pocc}
\end{equation}
If the number of strictly negative eigenvalues is not exactly \(N_{\rm occ}\),
the notebook falls back to filling the lowest \(N_{\rm occ}\) eigenvectors.

For the RG calculation we use the standard complex-fermion correlation matrix
\begin{equation}
  C_{ij}=\langle c_j^\dagger c_i\rangle,
  \qquad C=\Pocc.
  \label{eq:cdef}
\end{equation}
The equivalent \(\pm 1\) covariance convention used elsewhere in the repository is
\begin{equation}
  G = 2C^* - I.
\end{equation}
The entanglement calculations below are performed directly with \(C\), because
the eigenvalues of a restricted correlation matrix are the occupation numbers
entering the Gaussian entropy.

\section{Entropy from a restricted Gaussian state}

For any set \(A\) of single-particle orbitals, let \(R_A\) be the rectangular
restriction matrix that selects those orbitals.  The restricted correlation
matrix is
\begin{equation}
  C_A = R_A C R_A^\dagger.
  \label{eq:restricted-c}
\end{equation}
If \(\lambda_\ell\in[0,1]\) are the eigenvalues of \(C_A\), the von Neumann
entropy of the many-body Gaussian reduced density matrix is
\begin{equation}
  S(A)= -\sum_{\ell}
  \left[
    \lambda_\ell\ln\lambda_\ell
    +(1-\lambda_\ell)\ln(1-\lambda_\ell)
  \right].
  \label{eq:entropy}
\end{equation}
This follows from the Gaussian normal form
\begin{equation}
  \rho_A = Z_A^{-1}\exp\left(-\sum_\ell \xi_\ell f_\ell^\dagger f_\ell\right),
  \qquad
  \lambda_\ell = \frac{1}{1+e^{\xi_\ell}},
  \label{eq:rho-gaussian}
\end{equation}
where the \(f_\ell\) are natural fermion modes of the subsystem.

The diagnostic subsystem is a full-\(x\) strip of height \(A_y\),
\begin{equation}
  A(y_0,A_y)=
  [0,N_x)\times[y_0,y_0+A_y)
\end{equation}
with periodic wrapping in \(y\).  In the thermodynamic infrared limit expected
for two counter-propagating domain-wall channels, the averaged strip entropy
obeys
\begin{equation}
  \overline{S}(A_y)
  = m \log\left[
      \frac{N_y}{\pi}\sin\left(\frac{\pi A_y}{N_y}\right)
    \right]+b,
  \qquad
  c_{\rm est}=3m.
  \label{eq:logchord}
\end{equation}
For the full-\(x\) strip in the physical domain-wall benchmark, the target is
\begin{equation}
  c_{\rm est}\to c_{\rm eff}=1.
  \label{eq:target-c}
\end{equation}
This value is not the chiral central charge of a single edge.  It is the
Calabrese-Cardy effective central charge seen by a quasi-one-dimensional cut
that intersects both domain-wall channels.

\section{Supersite representation}

The RG is performed along \(y\).  We combine all orbitals in a fixed row into a
single supersite.  A microscopic orbital index is written as
\begin{equation}
  i=(y,\mu),
  \qquad
  y=0,\ldots,N_y-1,
  \qquad
  \mu=1,\ldots,q,
  \qquad
  q=2N_x.
  \label{eq:supersite-index}
\end{equation}
In the code, \(\mu\) is the combined \((x,\sigma)\) index.  With this ordering
the correlation matrix is a block matrix
\begin{equation}
  C_{y\mu,y'\nu}.
\end{equation}
At RG layer \(\ell\), the number of supersites is \(N_y^{(\ell)}\), while the
number of retained orbitals per supersite is kept fixed at \(q=2N_x\).  The
default sequence is therefore
\begin{equation}
  N_y^{(0)}=40,\qquad
  N_y^{(1)}=20,\qquad
  N_y^{(2)}=10.
  \label{eq:ny-flow}
\end{equation}

\section{Block natural-orbital RG}

A binary RG layer groups neighboring supersites,
\begin{equation}
  B_j=\{2j,2j+1\},
  \qquad
  j=0,\ldots,N_y^{(\ell)}/2-1.
\end{equation}
The corresponding two-row block contains \(2q\) fermion modes.  Let \(R_j\) be
the \(2q\times qN_y^{(\ell)}\) restriction matrix selecting those modes.  The
local block correlation matrix is
\begin{equation}
  C_{B_j} = R_j C^{(\ell)} R_j^\dagger.
  \label{eq:block-c}
\end{equation}
It is diagonalized as
\begin{equation}
  C_{B_j}=V_j n_j V_j^\dagger,
  \qquad
  n_j=\diag(n_{j,1},\ldots,n_{j,2q}).
  \label{eq:natural-orbitals}
\end{equation}
The columns of \(V_j\) define natural orbitals
\begin{equation}
  d_{j,a}=\sum_{\beta=1}^{2q}(V_j)_{\beta a}^\dagger c_{B_j,\beta},
\end{equation}
with occupations
\begin{equation}
  \langle d_{j,a}^\dagger d_{j,a}\rangle=n_{j,a}.
\end{equation}
For a Gaussian state, \(n_{j,a}\) has a direct entanglement interpretation.  If
\(n_{j,a}\) is close to \(0\) or \(1\), then that natural mode is nearly pure
inside the block and weakly entangled with the complement.  If
\(n_{j,a}\approx 1/2\), the mode is strongly mixed and carries long-distance
entanglement.

The notebook ranks modes using
\begin{equation}
  \chi_{j,a}=n_{j,a}(1-n_{j,a}).
  \label{eq:chi}
\end{equation}
This quantity is maximal at \(n=1/2\) and vanishes at \(n=0,1\).  In each
two-row block, the \(q\) modes with largest \(\chi_{j,a}\) are retained.  If
\(K_j\) denotes this set of retained natural-orbital labels, the local isometry
is
\begin{equation}
  W_j =
  \begin{pmatrix}
    | & & | \\
    v_{j,a_1} & \cdots & v_{j,a_q} \\
    | & & |
  \end{pmatrix},
  \qquad
  a_1,\ldots,a_q\in K_j,
  \label{eq:local-isometry}
\end{equation}
where \(v_{j,a}\) is the \(a\)-th column of \(V_j\).  Thus
\begin{equation}
  W_j^\dagger W_j=I_q.
\end{equation}
The global isometry is the block direct sum
\begin{equation}
  W^{(\ell)}=\bigoplus_j W_j.
  \label{eq:global-isometry}
\end{equation}
It maps the fine layer's \(qN_y^{(\ell)}\) modes to the coarse layer's
\(qN_y^{(\ell+1)}\) modes, with
\begin{equation}
  N_y^{(\ell+1)}=N_y^{(\ell)}/2.
\end{equation}
The coarse correlation matrix is
\begin{equation}
  C^{(\ell+1)}
  =
  \left(W^{(\ell)}\right)^\dagger
  C^{(\ell)}
  W^{(\ell)}.
  \label{eq:coarse-c}
\end{equation}
This is the central update of the notebook.

It is useful to also keep track of the embedding of each coarse mode into the
original microscopic Hilbert space.  Define
\begin{equation}
  E^{(0)}=I,
\end{equation}
and recursively
\begin{equation}
  E^{(\ell+1)}=E^{(\ell)}W^{(\ell)}.
  \label{eq:embedding}
\end{equation}
The \(x\)-profile plotted in the notebook is
\begin{equation}
  p^{(\ell)}(x)
  =
  \frac{
    \sum_{y,\sigma,a}
    \left|E^{(\ell)}_{(x,y,\sigma),a}\right|^2
  }{
    \sum_{x',y,\sigma,a}
    \left|E^{(\ell)}_{(x',y,\sigma),a}\right|^2
  }.
  \label{eq:x-profile}
\end{equation}
If the retained subspace is dominated by the domain-wall edge channels, this
profile should place enhanced weight near the two domain-wall positions.

\section{Relation to MERA}

The update in Eq.~\eqref{eq:coarse-c} has the same algebraic form as a Gaussian
MERA coarse graining after setting all disentanglers to the identity:
\begin{equation}
  C^{(\ell+1)}
  =
  \left(W^{(\ell)}\right)^\dagger
  \left(U^{(\ell)} C^{(\ell)} U^{(\ell)\dagger}\right)
  W^{(\ell)},
  \qquad
  U^{(\ell)}=I.
  \label{eq:mera-with-u}
\end{equation}
A genuine scale-invariant Gaussian MERA would go further.  One would optimize
local Gaussian unitaries \(U\) across block boundaries and local isometries
\(W\), then reuse the same tensors at every layer, at least after microscopic
transients.  In such a construction \(U\) and \(W\) are variational parameters
chosen to minimize the energy, maximize local disentangling, or impose an
approximate fixed-point condition on the coarse-grained correlation matrix.

The present RG deliberately avoids that harder variational problem.  Its
limitations are:
\begin{enumerate}
\item no boundary disentangler is optimized;
\item the isometries \(W_j\) are determined independently in each block;
\item the isometries are layer dependent;
\item the procedure is diagnostic rather than a compact ansatz for the whole
      many-body state.
\end{enumerate}
For this reason the notebook should be interpreted as a Gaussian entanglement
RG, not as a completed MERA construction.

\section{Scale-invariant Gaussian MERA prototype}

The companion notebook \texttt{scale\_invariant\_gaussian\_mera.ipynb}
implements the first variational step beyond the smoke-test RG.  The goal is
not to optimize an unconstrained \(U(80)\times{\rm Stiefel}(80,40)\) ansatz,
but to test whether a single repeated Gaussian layer can behave as an
approximately scale-invariant coarse grainer for the domain-wall state.

\subsection{Repeated layer geometry}

At a fixed RG layer, write the quasi-one-dimensional state as \(C\) on
\(N_y\) supersites with internal dimension \(q=2N_x\).  A binary MERA layer has
two staggered operations.  First, a repeated boundary disentangler
\(u\in U(2q)\) acts on the pairs
\begin{equation}
  (1,2),\;(3,4),\ldots,\;(N_y-1,0),
  \label{eq:boundary-pairs}
\end{equation}
where the final pair uses periodic wrapping.  The corresponding global unitary
is
\begin{equation}
  U = \bigoplus_{j=0}^{N_y/2-1} u_{(2j+1,\,2j+2)}.
  \label{eq:global-disentangler}
\end{equation}
Second, a repeated isometry \(w\in\mathbb{C}^{2q\times q}\) with
\begin{equation}
  w^\dagger w=I_q
  \label{eq:w-isometry}
\end{equation}
acts on the block pairs
\begin{equation}
  (0,1),\;(2,3),\ldots,\;(N_y-2,N_y-1).
  \label{eq:block-pairs}
\end{equation}
The global isometry is
\begin{equation}
  W = \bigoplus_{j=0}^{N_y/2-1} w_{(2j,\,2j+1)}.
  \label{eq:global-repeated-isometry}
\end{equation}
The single-particle correlation matrix is then updated by
\begin{equation}
  C' = W^\dagger U C U^\dagger W.
  \label{eq:si-mera-update}
\end{equation}
Equation~\eqref{eq:si-mera-update} is the finite-correlation-matrix form of a
Gaussian MERA layer.  Its difference from Eq.~\eqref{eq:coarse-c} is not the
algebraic form, but the variational restriction: the same \(u\) and \(w\) are
reused at every block and, in the scale-invariant test, at every layer.

\subsection{Initialization of the repeated isometry}

The repeated isometry is initialized from the translation-averaged two-row
correlation matrix on the block pairing in Eq.~\eqref{eq:block-pairs}.  Define
\begin{equation}
  \overline{C}_{B}
  =
  \frac{2}{N_y}
  \sum_{j=0}^{N_y/2-1}
  R_{2j,2j+1} C R_{2j,2j+1}^\dagger,
  \label{eq:avg-block-c}
\end{equation}
where \(R_{a,b}\) restricts to the two supersites \(a,b\).  Diagonalize
\begin{equation}
  \overline{C}_B = V_B n_B V_B^\dagger,
\end{equation}
rank the \(2q\) natural orbitals by
\begin{equation}
  \chi_a = n_a(1-n_a),
\end{equation}
and form \(w\) from the \(q\) columns of \(V_B\) with largest \(\chi_a\).
Thus the repeated isometry is the translationally averaged analogue of the
block-dependent natural-orbital map used in the smoke-test RG.

\subsection{Fixed-point covariance objective}

To test approximate scale invariance, the notebook compares translation-averaged
block correlators before and after coarse graining.  For the fine state define
\begin{equation}
  \overline{C}_{0r}
  =
  \frac{1}{N_y}
  \sum_{y=0}^{N_y-1}
  R_y C R_{y+r}^\dagger,
  \label{eq:avg-correlator}
\end{equation}
where \(R_y\) restricts to the \(q\) orbitals on supersite \(y\).  For the
coarse state \(C'\), the analogous object is \(\overline{C}'_{0r}\).  A
scale-invariant layer should map coarse separation \(r\) to fine separation
\(2r\), so the loss is
\begin{equation}
  \mathcal{L}
  =
  \sum_{r=0}^{r_{\max}}
  \omega_r
  \frac{
    \left\|
      \overline{C}'_{0r}-\overline{C}_{0,2r}
    \right\|_F^2
  }{
    \left\|
      \overline{C}_{0,2r}
    \right\|_F^2+\epsilon
  }.
  \label{eq:fixed-point-loss}
\end{equation}
The implementation uses \(\omega_r=(1+r)^{-1}\).  This objective is
gauge-sensitive because the retained coarse modes define an internal basis.
Consequently the fixed-point loss is interpreted together with the physical
diagnostics: the log-chord entropy coefficient, the entanglement spectrum, and
the retained \(x\)-profile.

\subsection{Active-subspace disentangler}

The first optimized disentangler is restricted to a small active subspace of
the boundary pair.  Let
\begin{equation}
  \overline{C}_{\partial}
  =
  \frac{2}{N_y}
  \sum_{j=0}^{N_y/2-1}
  R_{2j+1,2j+2} C R_{2j+1,2j+2}^\dagger
  \label{eq:avg-boundary-c}
\end{equation}
be the averaged two-row correlation matrix on the disentangler pairing.  Its
natural orbitals with largest \(\chi=n(1-n)\) define an orthonormal matrix
\begin{equation}
  B\in\mathbb{C}^{2q\times m},\qquad B^\dagger B=I_m,
\end{equation}
with \(m=12\) in the default run and \(m=4\) in smoke tests.  A small unitary on
this subspace is parameterized as
\begin{equation}
  u_{\rm act}=\exp K,\qquad K^\dagger=-K.
  \label{eq:active-exp}
\end{equation}
It is embedded into the full two-supersite disentangler by
\begin{equation}
  u = I_{2q} + B\left(u_{\rm act}-I_m\right)B^\dagger.
  \label{eq:active-embed}
\end{equation}
The entries of the anti-Hermitian generator \(K\) are optimized with
\texttt{scipy.optimize.minimize} using the fixed-point loss in
Eq.~\eqref{eq:fixed-point-loss}.  This gives a controlled, low-dimensional
variational step toward a true Gaussian MERA without attempting a fragile
full-matrix optimization.

\subsection{Acceptance diagnostics}

The scale-invariant MERA prototype is considered successful only if several
diagnostics agree.  The fixed-point loss should decrease after optimizing the
active disentangler; the repeated-layer entropy should remain close to
\(c_{\rm eff}=1\); the retained subspace should remain localized near the two
domain walls in \(x\); and the half-chain entanglement spectrum near
\(\lambda=1/2\) should be preserved under repeated coarse graining.  A reduced
loss without these physical signatures would be treated as a gauge artifact, not
as evidence for an infrared MERA fixed point.

\section{Self-consistent MERA v2}

The first scale-invariant prototype is useful because it exposes a failure mode.
It applies a microscopic natural-orbital isometry repeatedly, but the output of
one layer is not guaranteed to be a pure Gaussian fixed-point state in the same
internal gauge.  If the discarded modes remain entangled with the retained
subspace, the raw compressed matrix
\begin{equation}
  C_{\rm raw}'=W^\dagger U C U^\dagger W
  \label{eq:raw-compression}
\end{equation}
has eigenvalues spread through \([0,1]\), rather than being close to a
projector.  Repeating this mixed compression can create spurious interval
entropy unrelated to the edge CFT.  This is precisely the pathology that the
second notebook version is designed to diagnose.

\subsection{Raw and purified coarse states}

The v2 notebook separates the raw compressed matrix from the pure Gaussian
teacher state used for repeated-layer diagnostics.  After Eq.~\eqref{eq:raw-compression},
one constructs
\begin{equation}
  C_{\rm pure}'=\Pi_{1/2}(C_{\rm raw}'),
  \label{eq:pure-projection}
\end{equation}
where \(\Pi_{1/2}\) denotes diagonalizing \(C_{\rm raw}'\) and occupying the
largest half of its eigenmodes.  Thus \(C_{\rm pure}'\) is the nearest
half-filled projector in the eigenbasis of \(C_{\rm raw}'\).  The raw matrix is
not discarded: the notebook reports the idempotency defect
\begin{equation}
  \delta_{\rm idem}(C_{\rm raw}')
  =
  \frac{\left\|C_{\rm raw}'{}^2-C_{\rm raw}'\right\|_F}
       {\sqrt{\dim C_{\rm raw}'}}.
  \label{eq:idem-defect}
\end{equation}
If this defect is large, the layer has not yet learned a true MERA
disentangler-isometry pair, even if the purified teacher state has stable
long-distance observables.

\subsection{Discarded-sector diagnostics}

Let \(d\in\mathbb{C}^{2q\times q}\) be an orthonormal complement to the retained
isometry \(w\), so that \((w,d)\) is unitary on a two-supersite block.  For the
post-disentangler local block correlation \(C_B^U\), define
\begin{align}
  C_{KK} &= w^\dagger C_B^U w, &
  C_{DD} &= d^\dagger C_B^U d, &
  C_{KD} &= w^\dagger C_B^U d .
\end{align}
The discarded-sector occupations are the eigenvalues \(n_{d,a}\) of \(C_{DD}\).
A successful layer should make these modes nearly pure and weakly coupled to
the retained sector:
\begin{equation}
  \overline{\chi}_{D}
  =
  \frac{1}{q}\sum_{a=1}^{q} n_{d,a}(1-n_{d,a})
  \approx 0,
  \qquad
  \|C_{KD}\|_F^2 \approx 0.
  \label{eq:discard-diagnostics}
\end{equation}
These quantities directly test the local MERA requirement that discarded modes
be disentangled, rather than merely removed.

\subsection{Self-consistent isometry update}

The v2 notebook also avoids blindly reusing the microscopic \(w\).  Starting
from the averaged natural-orbital isometry, it performs a teacher iteration:
\begin{enumerate}
\item apply one raw MERA layer using the current \(w\);
\item project the coarse matrix to \(C_{\rm pure}'\);
\item extract a new averaged natural-orbital isometry \(w_{\rm next}\) from
      \(C_{\rm pure}'\);
\item align \(w_{\rm next}\) to the old output gauge by a Procrustes/polar
      rotation;
\item damp-update \(w\) and re-orthonormalize.
\end{enumerate}
The gauge alignment solves
\begin{equation}
  R_*=\arg\min_{R^\dagger R=I_q}\|w_{\rm next}R-w\|_F,
  \label{eq:procrustes}
\end{equation}
implemented by the polar factor of \(w_{\rm next}^\dagger w\).  The damped
update is
\begin{equation}
  w \leftarrow {\rm qf}\!\left[(1-\eta)w+\eta\,w_{\rm next}R_*\right],
  \label{eq:damped-w}
\end{equation}
where \({\rm qf}\) denotes QR re-orthonormalization and \(\eta\) is the damping
parameter.

\subsection{Composite v2 objective}

The active-subspace disentangler in v2 is optimized against a composite loss
whose leading terms enforce local MERA consistency:
\begin{equation}
  \mathcal{L}_{\rm v2}
  =
  \lambda_{\rm idem}\delta_{\rm idem}^2
  +\lambda_{\rm disc}\overline{\chi}_{D}
  +\lambda_{\rm cross}\overline{\|C_{KD}\|_F^2}
  +\lambda_{\rm fp}\mathcal{L}_{\rm corr}.
  \label{eq:v2-loss}
\end{equation}
The default notebook weights are
\begin{equation}
  \lambda_{\rm idem}=1,\qquad
  \lambda_{\rm disc}=2,\qquad
  \lambda_{\rm cross}=1,\qquad
  \lambda_{\rm fp}=0.05.
\end{equation}
The small coefficient of \(\mathcal{L}_{\rm corr}\) reflects the fact that the
block-correlator fixed-point loss is gauge-sensitive.  For a first MERA smoke
test, purity and discarded-sector polarization are more direct indicators that
the layer is learning a legitimate coarse graining.

The v2 construction remains a controlled prototype.  It is not a full
unconstrained optimization over \(u\in U(2q)\) and
\(w\in{\rm Stiefel}(2q,q)\).  Its purpose is to determine whether the visible
domain-wall subspace can be made stable under repeated quasi-one-dimensional
Gaussian coarse graining before attempting a more expensive variational MERA.

\section{Finite-depth Gaussian MERA attempt}

The follow-up notebook
\texttt{finite\_depth\_gaussian\_mera.ipynb} removes the strongest assumption
in the scale-invariant prototypes: it does not force a single microscopic
tensor pair to act at every scale.  Instead, each layer has its own Gaussian
maps
\begin{equation}
  C^{(\ell+1)}
  =
  W_\ell^\dagger U_\ell C^{(\ell)} U_\ell^\dagger W_\ell,
  \qquad
  \ell=0,1,\ldots ,
  \label{eq:finite-depth-layer}
\end{equation}
with \(U_\ell=\oplus_B u_{\ell,B}\) acting on the odd-even block boundaries and
\(W_\ell=\oplus_B w_{\ell,B}\) acting on the even-odd blocks that are then
coarse grained.  In the translationally averaged version used in the notebook,
all \(u_{\ell,B}\) in a given layer are tied to the same active-subspace
unitary \(u_\ell\), and all \(w_{\ell,B}\) are tied to the same two-row
isometry \(w_\ell\).  The crucial difference from a scale-invariant ansatz is
that \(u_0,w_0\) and \(u_1,w_1\) are learned independently.

For a fixed layer, the notebook first applies a candidate boundary
disentangler,
\begin{equation}
  C_U^{(\ell)}=U_\ell C^{(\ell)}U_\ell^\dagger ,
\end{equation}
and then forms the averaged two-row block covariance
\begin{equation}
  \overline{C}_{B}^{(\ell)}
  =
  \frac{2}{N_y^{(\ell)}}\sum_{j=0}^{N_y^{(\ell)}/2-1}
  R_{2j,2j+1}C_U^{(\ell)}R_{2j,2j+1}^\dagger .
\end{equation}
The isometry \(w_\ell\) is extracted only after this disentangling step, by
diagonalizing \(\overline{C}_B^{(\ell)}\) and retaining the \(q\) orbitals with
largest \(\chi_a=n_a(1-n_a)\).  This ordering is closer to the MERA logic:
the disentangler should make the discarded sector locally pure before the
isometry chooses which modes to keep.

The active \(u_\ell\) search used here is deliberately modest.  An active basis
is chosen from the most mixed eigenvectors of the averaged boundary covariance,
and random anti-Hermitian generators \(K=-K^\dagger\) are exponentiated inside
that subspace,
\begin{equation}
  u_\ell
  =
  I_{2q}
  +B_\ell\left(e^K-I_m\right)B_\ell^\dagger .
\end{equation}
Each accepted trial lowers the local objective
\begin{equation}
  \mathcal{L}_{\rm fd}
  =
  \lambda_{\rm idem}\delta_{\rm idem}^2
  +\lambda_{\rm disc}\overline{\chi}_{D}
  +\lambda_{\rm cross}\overline{\|C_{KD}\|_F^2}
  +\lambda_{\rm fp}\mathcal{L}_{\rm corr}.
  \label{eq:finite-depth-loss}
\end{equation}
The objective is the same physical test as in Eq.~\eqref{eq:v2-loss}, but it
is now applied separately at each scale rather than to a repeated fixed-point
map.

The full \(N_x=20,N_y=40\) run gives an instructive outcome.  The independent
finite-depth natural-orbital flow with \(u_\ell=I\) remains nearly pure through
two layers, with raw idempotency defects of order \(10^{-5}\), and its fitted
central charge stays close to one:
\begin{equation}
  c_{\rm est}^{(0)}\simeq1.0000,\qquad
  c_{\rm est}^{(1)}\simeq1.0538,\qquad
  c_{\rm est}^{(2)}\simeq1.0162 .
\end{equation}
This is the first positive MERA-like result: the catastrophic layer-two
entropy blowup of the repeated scale-invariant ansatz is not present once the
isometry is allowed to change with scale.  The retained \(x\)-profile remains
peaked near the domain walls, and the discarded natural-orbital occupations
are strongly polarized.

The active random disentangler does not yet improve the raw finite-depth flow.
With the present small active search it slightly worsens the second-layer
entropy drift, while a purified-teacher variant gives \(c_{\rm est}\approx1\)
only because the state is explicitly projected back to a half-filled Gaussian
projector after each layer.  Thus the main lesson is not that the active
optimizer is successful, but that the scale-invariance constraint was the
dominant failure mode.  A practical next step is to use the learned finite-depth
\(\{w_\ell\}\) as data: gauge-align them across layers, compare their
difference, and only then fit a reusable scale-invariant tensor.

\subsection{Attempted scale-invariant fit from finite-depth tensors}

The notebook now performs this last diagnostic explicitly.  Let \(w_0\) be the
isometry learned for the \(40\to20\) step and \(w_1\) the isometry learned for
the \(20\to10\) step.  Since the output basis of each isometry is arbitrary,
the first comparison removes the right output gauge by solving
\begin{equation}
  R_*=\arg\min_{R^\dagger R=I_q}\|w_1R-w_0\|_F .
\end{equation}
Equivalently, if \(w_1^\dagger w_0=U\Sigma V^\dagger\), then
\(R_*=UV^\dagger\).  The notebook records
\begin{equation}
  \Delta_{\rm out}
  =
  \frac{\|w_1R_*-w_0\|_F}{\|w_0\|_F},
  \qquad
  \Delta_P
  =
  \frac{\|w_1w_1^\dagger-w_0w_0^\dagger\|_F}{\sqrt q},
\end{equation}
together with the principal singular values of \(w_0^\dagger w_1\).

For the full run the layer tensors are not close:
\begin{equation}
  \Delta_{\rm out}\simeq0.875,\qquad
  \Delta_P\simeq0.957
  \qquad
  \text{for the } u=I \text{ finite-depth flow}.
\end{equation}
Even the active and purified variants only reduce these numbers to roughly
\(\Delta_{\rm out}\simeq0.77\)--\(0.78\) and
\(\Delta_P\simeq0.85\)--\(0.86\).  This is too large to interpret the two
finite-depth layers as small gauge deformations of a single scale-invariant
isometry.

As a direct stress test, the notebook constructs an averaged candidate
\begin{equation}
  w_* = {\rm qf}\!\left[\frac{1}{2}(w_0+w_1R_*)\right],
\end{equation}
and reuses \(w_*\) as a repeated scale-invariant tensor.  This fails
immediately: for the \(u=I\) finite-depth tensors the repeated fit gives
\(c_{\rm est}\sim 3.8\times10^2\) after one layer and has raw idempotency
defect of order \(10^{-1}\).  Thus the finite-depth construction is a positive
real-space RG, but it is not yet evidence for an ordinary one-site-periodic
scale-invariant MERA.  A viable scale-invariant construction would need either
a better coarse-site gauge, a nontrivial repeating disentangler optimized
together with \(w\), or a larger unit cell in scale.

\subsection{Canonical gauge, period-two test, and reduced wall channels}

The finite-depth notebook was extended with three further tests of this
obstruction.  First, each coarse state is put in a canonical onsite gauge.
For a layer-\(\ell\) covariance, define the translation-averaged onsite and
nearest-neighbor matrices
\begin{equation}
  \overline C_0^{(\ell)}
  =
  \frac{1}{N_y^{(\ell)}}\sum_y C_{y,y}^{(\ell)},
  \qquad
  \overline C_1^{(\ell)}
  =
  \frac{1}{N_y^{(\ell)}}\sum_y C_{y,y+1}^{(\ell)} .
\end{equation}
The unitary \(G_\ell\) is chosen to diagonalize \(\overline C_0^{(\ell)}\),
with residual near-degenerate subspaces fixed by the Hermitian part of
\(\overline C_1^{(\ell)}\).  The layer isometry is then compared in the
gauge-fixed form
\begin{equation}
  \widetilde w_\ell
  =
  (G_\ell^\dagger\oplus G_\ell^\dagger)\,
  w_\ell\,G_{\ell+1}.
\end{equation}
This asks whether the apparent scale dependence was merely a poor choice of
coarse orbital basis.

It was not.  For the \(N_y=40\) run the canonical-gauge mismatch between the
two finite-depth layers is still order one:
\begin{equation}
  \Delta_{\rm out}\simeq0.878,\qquad
  \Delta_P\simeq0.957
  \qquad (u=I).
\end{equation}
The active and purified variants only reduce this to
\(\Delta_{\rm out}\simeq0.78\)--\(0.80\).  Thus canonical onsite gauge fixing
does not turn the finite-depth tensors into a one-tensor fixed point.

Second, an optional \texttt{GENTRG\_DEEP\_NY80=1} run performs
\(80\to40\to20\to10\), with active disentangler optimization disabled so that
the deeper finite-depth natural-orbital flow is affordable.  The entropy flow
remains good: \(c_{\rm est}\simeq1.113,1.059,1.027\) after the three coarse
layers, with raw idempotency defects of order \(10^{-5}\).  However, the
period-two hypothesis also fails.  Comparing layer \(0\) and layer \(2\) in
canonical gauge gives
\begin{equation}
  \Delta_{\rm out}^{(0,2)}\simeq0.869,\qquad
  \Delta_P^{(0,2)}\simeq0.940 .
\end{equation}
The finite-depth flow is therefore not a simple period-two scale-invariant
cycle in this row-orbital basis.

Third, the notebook tests a deliberately simple domain-wall-channel projection:
keep only the microscopic orbitals with \(x\) within one lattice spacing of
each domain wall, giving \(q_{\rm eff}=12\), purify the projected covariance to
a half-filled Gaussian projector, and run the same natural-orbital finite-depth
RG.  This projection is too crude.  The initial projected state already has
\(c_{\rm est}\simeq2\), and under coarse graining it drifts to values between
roughly \(3\) and \(5\).  Although the retained subspace is by construction
wall-localized, the simple real-space wall window does not isolate the single
low-energy Dirac channel.  A useful reduced theory must be built from
Hamiltonian or correlation eigenchannels localized at the walls, not from a
hard \(x\)-space window.

\section{Why this is a useful smoke test}

The physical question is whether the flattened-Hamiltonian ground state has an
infrared decomposition into gapped two-dimensional bulk modes plus gapless
domain-wall modes.  The smoke-test logic is:
\begin{align}
  \text{gapped local bulk mode}
  &\Longrightarrow n\approx0 \text{ or }1
  \Longrightarrow \chi=n(1-n)\approx0,
  \\
  \text{IR entangled edge mode}
  &\Longrightarrow n\approx1/2
  \Longrightarrow \chi\approx1/4.
\end{align}
Thus a successful local Gaussian mode sort should discard polarized modes and
retain mixed modes.  The retained subspace should then be localized near the
domain walls in \(x\), while the strip entropy should continue to fit the
periodic log-chord form with \(c_{\rm est}\) close to one.

Failure is also informative.  If the retained modes are not wall localized, the
full-\(x\) entropy slope rapidly drifts, or the discarded modes are not
polarized, then an optimized scale-invariant MERA is premature.  The likely
fixes would be to include explicit disentanglers, enlarge the local block with
buffer rows, first project onto an edge-mode subspace, or increase the system
size to reduce finite-size mixing between the two domain walls.

\section{Notebook observables}

The notebook records four primary diagnostics at each layer:
\begin{enumerate}
\item the averaged full-\(x\) strip entropy and the fit
      \(c_{\rm est}=3m\);
\item the kept and discarded natural-orbital occupation spectra;
\item the retained-subspace \(x\)-profile \(p^{(\ell)}(x)\);
\item the half-chain entanglement spectrum near \(\lambda=1/2\).
\end{enumerate}
Together these observables test the minimal necessary ingredients for a future
scale-invariant Gaussian MERA: local mode polarization, wall localization of
the retained modes, and persistence of the expected \(c_{\rm eff}=1\) infrared
entanglement coefficient.

\end{document}
